% 原书第 421 页（PDF 第 434 页）
\BookPage{421}{434}
\BookHeading{记一件康熙朝玻璃器的最高成就}{记一件康熙朝玻璃工艺的最高成就\\[3pt]—玻璃胎画珐瑯牡丹蓝地胆瓶*}{嵇若昕**}
\BookSubheading{楔子}

2010年上半年，于北京故宫博物院（以下简称「北京故宫」）进行为期五个月的考察与研究，主要是对于国立故宫博物院（以下简称「台北故宫」）的藏品溯源。众人皆知，台北故宫的藏品主要是原北平故宫博物院南迁文物中精品，本次研究计划乃希望能一窥台北故宫藏品在清季贮存于紫禁城时的大致状况。因此查阅北京故宫图书馆所收藏的内廷陈设档案，遂成为刻不容缓的工作。

北京故宫所藏内廷陈设档案有数百种，居停期间就台北故宫藏品内容与个人研究偏好，查阅了近百种的陈设档，最早的有康熙三十三年（1694）正月建立的《（坤宁宫）陈设帐》，最晚的是宣统十六年（民国十三年，1924）八月（农历）建立的《王、罗大人提出古铜册》，同年阳历十一月五日末代皇帝溥仪就被请出宫了。在个人查阅的这批陈设档中，以俗称「多宝格」的清宫百什件最多，¹但是数本有关清代珐瑯器的陈设档案，则充分展现了清代珐瑯器历经康、雍、干三帝的关心，内廷从摸索、研制以致登峰造极的历程，以及清中后期紫禁城内对其重视的现象。

\vfill
\begin{BookNotes}
\BookNote{}{2010年上半年，个人承香港「利荣森纪念交流计划赞助，前往北京故宫进行考察与研究，收获甚丰。本文乃本项赞助下的第一篇心得报告，特此致谢。}
\BookNote{**}{作者任职于国立故宫博物院。}
\BookNote{1.}{关于清宫中的百什件，详见拙文，〈上下五千年，东西十万里—清宫中的百什件〉，原载于《故宫文物月刊》，294期，2007年9月，页4-15，经局部修订后收录于《双溪文物随笔》（台北：故宫，2011），页197-207。}
\end{BookNotes}

% 原书第 422 页（PDF 第 435 页）
\BookPage{422}{435}
\noindent 乾清宫珐瑯器陈设档

北京故宫图书馆珍藏的专载清代康、雍、干三朝珐瑯器档案有道光十五年(1835）七月初一日建立的《珐瑯玻璃宜兴磁胎陈设档案》（封面绿色号签为「故宫博物院文献馆、登记:设8（墨书）」），2与《珐瑯金银铜胎陈设档案》（封面绿色号签为「故宫博物院文献馆、登记:设7（墨书）」）；光绪元年（1875）十一月初七日建立的《珐瑯玻璃宜兴磁胎陈设档案》（封面绿色号签为「故宫博物院文献馆、登记：设322（原墨书363）」，封底红色号签为「故宫博物院文献馆、册字26847（蓝色打印后又用墨笔涂销）」），与《珐瑯金银铜胎陈设档案》（封面绿色号签为「故宫博物院文献馆、登记:设321（墨书）」，封底红色号签为「故宫博物院文献馆、册字26848（蓝色打印后又用墨笔涂销）」）；光绪二年（1876）三月初四日建立的《金银铜宜兴瓷胎玻璃珐瑯档案》（封面绿色号签为「故宫博物院文献馆、登记：设477（墨书）」，封底红色号签为「故宫博物院文献馆、口（漫澳不清）字311?3（蓝色打印）」）；以及光绪二十八年（1902）二月初二日建立的《明殿现设金银铜宜兴瓷玻璃瑯档》（封面红色号签为「故宫博物院文献馆、登记:设415（墨书）」、口（漫澳不清）字29582（蓝色打印）」）。其中道光十五年七月初一日建立的《瑯玻璃宜兴磁胎陈设档案》、《珐瑯金银铜胎陈设档案》与光绪二年三月初四日建立的《金银铜宜兴瓷胎玻璃珐瑯档案》几乎每页都骑缝钤盖了汉满朱文「广储司印」。

经详细比对，这些档册应是当年内廷盘点后所制作的文物帐册，例如道光十五年七月初一日建立的《珐瑯金银铜胎陈设档案》，与光绪元年十一月初七日建立的《珐瑯金银铜胎陈设档案》内容与条目次序基本相同，载录顺序是先依据质材，再依据朝代，首载康熙朝的金胎珐瑯器，依序是雍正朝、乾隆朝的金胎珐瑯器，接着是康雍干三朝的铜胎珐瑯器，倒数第三页最后两行又开始补录乾隆朝金胎珐瑯器；扣除封面与封底，两册内文皆各仅二十页，即使两者第十九页第一行：「金胎珐瑯

\vfill
\begin{BookNotes}
\BookNote{2.}{廖宝秀，〈乾隆磁胎洋彩综述〉，收于《华丽彩瓷—乾隆洋彩》（台北：故宫，2008年），页13，刊登了这份档案的封面与第八页至第十一页的内页图像。廖宝秀，〈茶器瑰宝一康熙宜兴胎画珐瑯折枝花卉〉，《故宫文物月刊》，302期，2008年5月，页4，与《茶韵茗事一故宫茶话》（台北：故宫，2011)，页102，皆刊登了这份档案的第一页与第二页的内页图像。}
\end{BookNotes}

% 原书第 423 页（PDF 第 436 页）
\BookPage{423}{436}
\noindent 盖豆壹对」下方浮贴二张黄色签条，一张上墨书「咸丰三年六月十八日小太监平顺传上要去金胎瑯盖豆一件」，另一张墨书:「现存一件」，也皆相同。仅只有道光朝者第四页第一条：「铜胎珐瑯盘壹件、进瑯补蜡」之「盘」字乃用黄色签另书写后黏贴上，可看出下方原为「花篮」二字；光绪朝者则径墨书「铜胎珐瑯盘壹件、进瑯补蜡」。

道光十五年七月初一日建立的《珐瑯玻璃宜兴磁胎陈设档案》与光绪元年十一月初七日建立的《珐瑯玻璃宜兴磁胎陈设档案》也是基本相同，载录顺序基本上也是先依据质材，再依据朝代，首载仅一件的康熙朝的玻璃胎珐瑯器，接着是乾隆朝的玻璃胎珐瑯器，然后是康熙朝宜兴胎珐瑯器，康熙朝瓷胎珐瑯器随其后载录；雍正朝与乾隆朝没有宜兴胎珐瑯器，直接载录瓷胎珐瑯器。扣除封面与封底，道光十五年与光绪元年建立的宜兴胎、玻璃胎与瓷胎珐瑯器档案皆各二十七页。但是例如前者第一页第十三行:「宜兴胎画珐瑯包袱式壶壹件」的下方黏贴黄签:「咸丰四年十月初五日小太监日传上要去，赏皇后用宜兴胎画珐瑯包袱式壶一件」；第十七行：「宜兴胎画珐瑯小提梁壶」，「小」字乃在「珐瑯」二字旁边贴黄签墨书而成；最后一行：「磁胎画珐瑯五彩西番花红地茶碗肆对」下方黏贴黄签：「现存一匣」。后者第一页则径删去「宜兴胎画珐瑯包袱式壶壹件」条，「宜兴胎画珐瑯小提梁壶」的「小」字直接书写。3

至于光绪二年三月初四日建立的《金银铜宜兴瓷胎玻璃珐瑯档案》与光绪二十八年二月初二日建立的《明殿现设金银铜宜兴瓷玻璃珐瑯档》，前者内文有四十七页，后者内文高达八十七页。后者这八十七页中，其第一页至第四十页载录当时乾清宫明殿陈设的宝座、图书（如《国朝宫史》、《皇朝礼器图》等)、器用等，另外还包括「年节陈设」的各式盆景、「如是宝明间布堂安供」的佛像与法器、「里间面东宝座床」附近陈设的器用（如各式佛像、朝珠等）；第四十一页即开始载录光绪二十八年二月初二日总管李长喜等陆续查得乾清宫库收三朝珐瑯」器，包含金胎、银胎、铜胎、瓷胎、宜兴胎、玻璃胎者，有掐丝珐瑯、画珐瑯、广珐瑯、西洋珐瑯之分，也有径称珐瑯者。经详细比对这两份珐瑯档案，后者自第

\vfill
\begin{BookNotes}
\BookNote{3.}{廖宝秀，〈茶器瑰宝—康熙宜兴胎画珐瑯折枝花卉〉，原载于《故宫文物月刊》，302期，2008年5月，页4-19，经局部修订后收录于《茶韵茗事—故宫茶话》，页101-115。}
\end{BookNotes}

% 原书第 424 页（PDF 第 437 页）
\BookPage{424}{437}
\noindent 四十一页起乃抄录前者而有所损益，例如前者若原书写错误，而用黄签另墨书正确字黏贴遮掩者，后者皆径书写正确文字；前者若原多写了字，用黄纸遮蔽，后者则径自不再书写；前者若已有黄签说明皇帝下令提取瑯器，后者或径自经统计仅书剩下的件数；若整个条目的珐瑯器都提取无剩，后者便径自删除不再书写此条目。（详见附表）至于载录的次序，与前述四册不同，乃先依据朝代，再依据质材，因此首载康熙朝金胎珐瑯器，次为康熙朝铜胎珐瑯器，三为康熙朝宜兴胎珐瑯器，四乃康熙朝瓷胎珐瑯器，五即仅一件的康熙朝玻璃胎珐瑯器；雍正朝者依序是金胎，次为铜胎，然后是瓷胎珐瑯器；乾隆朝者除了没有宜兴胎者外，余与康熙朝之次序相同，乾隆朝玻璃胎珐瑯器殿后。

从前述的比对中可知，这些珐瑯器原贮于乾清宫的文物库房，六本陈设档虽然有两本载录金胎、银胎、铜胎珐瑯器，两本载录瓷胎、玻璃胎、宜兴胎珐瑯器，另两本则合并两者。但是比对四本的金属胎（金、银、铜胎）与其他胎质（玻璃、宜兴、瓷胎）珐瑯器的内容，基本相同。因此，各类胎质珐瑯器在道光十五年七月初一日完成清点建立帐册后，光绪元年十一月初七日再度清点后又建立一本当时的帐册，其间珐瑯器除皇帝传旨赏赐外，并无损益。

上个世纪八年代初，北京故宫的朱家滘在其发表的〈清代画珐瑯制器考一《工艺美术史料汇编》之一〉文末，「附录」中记载一本道光十五年七月一日建立的《乾清宫珐瑯、玻璃、宜兴磁胎陈设档》，「每页有『内务府广储司』骑缝印。」4，即是本文所述道光十五年（1835）七月初一日建立的《珐瑯玻璃宜兴磁胎陈设档案》，档册封面书名签条并无「乾清宫」三字，且骑缝印之印文乃汉满文朱文「广储司印」

经比对民国十四年（1925）三月二十日刊行的《清宫物品点查报告·第一编·第三册·端凝殿左右各屋》内容，可知前述这些珐瑯器，于溥仪出宫前后大多收贮于端凝殿左右各屋。端凝殿在乾清宫东庑，其南、北侧各三屋，南侧文物有整

\vfill
\begin{BookNotes}
\BookNote{4.}{朱家溍，〈清代画珐瑯制器考一《工艺美术史料汇编》之一〉，《故宫博物院院刊》，1982年3期，页74-76与96。道光朝内廷述及瓷器时，用「磁」字，朱家溍在此文中即如此使用。}
\end{BookNotes}

% 原书第 425 页（PDF 第 438 页）
\BookPage{425}{438}
\noindent 批康雍干三朝松花石砚，北侧小库或即收贮三朝珐瑯器（见下文）。由此可知，所谓乾清宫的库房应是端凝殿南北两侧小屋。此外，《清宫物品点查报告·第一编·第三册·端凝殿左右各屋》书中还记载着当时清室善后委员会点查文物时，在端凝殿两侧屋宇内清点出的《乾清宫东西暖阁陈设档案》（点查号：列五三三，一𬘡拾册）、《坤宁宫陈设档》（点查号：列五三四，一𬘡四册）、《仿古、时做上次等陈设档案》（点查号：列五三五，一継五册）、《珐瑯各种胎陈设档案》（点查号：列五三六，一𬘡二册）、《续入库贮陈设油露药材等项档案》（点查号:列五三七，一册）、《乾清宫东西暖阁档案》（点查号:列五三八，一四本）、《坤宁宫东暖阁档案》（点查号：列五三九，一𬘡二册、光绪二年三月）、《仿古、时做上次等玉器陈设档案》（点查号：列五四，一𬘡二册、光绪二年三月）、《续入库收油露药材档案》（点查号:列五四一，一册、光绪二年三月）、《金银铜宜兴瓷胎玻璃珐瑯档案》（点查号：列五四二，一册、光绪二年三月）、《明殿现设金银铜宜兴瓷玻璃珐瑯档案》（点查号：列五四三，一册、光绪二十八年二月），在最后这册的条目旁还注明：「以上十一号档案提存事务室」。也就是说，前述六本清内廷珐瑯器陈设档，有两册（光绪二年三月初四日建立的《金银铜宜兴瓷胎玻璃珐瑯档案》与光绪二十八年二月初二日建立的《明殿现设金银铜宜兴瓷玻璃珐瑯档案》）确定在清季民初时收贮于端凝殿两侧小屋中，至于另两册《珐瑯各种胎陈设档案》（点查号：列五三六）究竟是金属胎珐瑯器陈设档？还是玻璃胎、宜兴胎与瓷胎珐瑯器的陈设档，俟考！至于是道光十五年建立者？抑或光绪元年建立者？更不可知。

依据前述清宫瑯器陈设档案所载，当时乾清宫库房端凝殿屋宇收贮的康雍干三朝各种胎质的珐瑯器有千余件之多，瓷胎画珐瑯器就有四百件左右（不含洋彩器），以康熙朝宜兴胎珐瑯器为例，道光十五年的陈设档案所载录者有二十件，光绪朝的档册则减一件。然而无论是那一份档册，康熙朝玻璃胎珐瑯器都仅有一件，

\vfill
\begin{BookNotes}
\BookNote{5.}{吴振棫，《养吉斋丛录》，转引自章唐容辑，《清宫述闻》（台北：文海，1969)，卷4，页317，载：「端凝殿南三楹藏康熙、雍正、乾隆所用砚墨，其砚悉以松花江石为之，三朝各四十枚，形式不一。」}
\BookNote{6.}{清室善后委员会刊行，《清宫物品点查报告·第一编·第三册·端凝殿左右各屋》（北平：故宫，1925），卷四，页203下-204上。}
\end{BookNotes}

% 原书第 426 页（PDF 第 439 页）
\BookPage{426}{439}
\noindent 即「玻璃胎珐瑯牡丹蓝地胆瓶」。

康熙朝玻璃胎画珐瑯牡丹蓝地胆瓶

朱家滘在其〈清代画珐瑯制器考一《工艺美术史料汇编》之一〉文章之附录中刊录了原藏紫禁城端凝殿「左右各屋」的瓷胎画珐瑯器，说明：所记文物品物乃「照录原盛匣盖上的名称」，康熙朝的「玻璃胎画珐瑯牡丹蓝地胆瓶」为「康熙年制」，7清室善后委员会的点查报告作「玻璃胎画珐瑯牡丹蓝地胆瓶」，且「带盖无款」，点查号为列三六738；但是清朝道光与光绪年间的乾清宫珐瑯器陈设档中皆作「玻璃胎珐瑯牡丹蓝地胆瓶」，但列于「康熙款玻璃胎」类目中。9

这件点查号为列三六73的玻璃胎胆瓶，为台北故宫藏品，编号为故瓷17588，民国十四年三月五日上午经清室善后委员会出组人员（吴钟麟、陈任中、窦传照、佟金德）点查，1°曾长期深埋外双溪山洞库房。虽然民国七十四年(1985）这件并无款识的画瑯玻璃器已公布于世人面前，作者张临生（时任器物处处长）于文中依据此器器形、瑯纹饰工艺，认为「十分可能是康熙时代的作品」，11但是似乎被西方学者误以为是一件鼻烟壶。12民国九十五年（2006），这件康熙朝画珐瑯玻璃器于「盛世工艺一清康熙、雍正、乾隆」展览中展出，但仅引起少数关心康熙朝玻璃器的学者注意。翌年（2007），当时任职台北故宫的施静菲于

\vfill
\begin{BookNotes}
\BookNote{7.}{朱家溍，〈清代画珐瑯制器考一《工艺美术史料汇编》之一〉，页74。}
\BookNote{8.}{清室善后委员会刊行，《清宫物品点查报告·第一编·第三册·端凝殿左右各屋》，卷四，页186上。}
\BookNote{9.}{道光十五年七月初一日建立的《珐瑯玻璃宜兴磁胎陈设档案》与光绪元年十一月初七日建立的《法瑯玻璃宜兴磁胎陈设档案》皆列于第1页，前者图见廖宝秀，〈茶器瑰宝一康熙宜兴胎画珐瑯折枝花卉〉，《故宫文物月刊》，302期，2008年5月，页4、图一，与《茶韵茗事一故宫茶话》，页102、图一。光绪二年三月初四日建立的《金银铜宜兴瓷胎玻璃珐瑯档案》，列于页7；光绪二十八年二月初二日建立的《明殿现设金银铜宜兴瓷玻璃法瑯档》，列于页47。}
\BookNote{10.}{清室善后委员会刊行，《清宫物品点查报告·第二辑》（北京：线装书局景印），页2-461，「点查端凝殿左右屋人员统计表。}
\BookNote{11.}{张临生，〈清朝的玻璃料器〉，《故宫文物月刊》，3卷9期（总33)，页4-13。Chang Lin-sheng, "Ch'ing Dynasty Imperial Enamelled Glassware," Arts of Asia, Vol. 21, No.3 (May-June, 1991), pp.95-107.}
\BookNote{12.}{Peter Hardie, "Glass in China: Questions and Some Answers," Vormen Uit Vuur, No. 158(1996 / 2), pp. 2-14.}
\end{BookNotes}

% 原书第 427 页（PDF 第 440 页）
\BookPage{427}{440}
\noindent 其〈十八世纪东西交流的见证一清宫画珐瑯工艺在康熙朝的建立〉文中再度讨论这件玻璃器，名为「玻璃胎画珐瑯蓝地花卉瓶」；13第三年，台北故宫退休同仁廖宝秀在其〈茶器瑰宝一康熙宜兴胎画珐瑯折枝花卉〉一文中首度比对道光十五年七月一日建立的《珐瑯玻璃宜兴磁胎陈设档案》的记载，附带公布这件玻璃器新摄数位影像彩图，并指出:「此为《陈设档》上记载的唯一康熙玻璃胎画珐瑯。」14其实这件玻璃器原贮于一木匣中，匣盖上标示为「康熙年制」。（见下文，图一）

\BookArt{p440-01.jpg}{0.952}{0.314}

上个世纪，器物学界多将瓷胎画珐瑯器称为珐瑯彩瓷或径称珐瑯彩，玻璃胎画珐瑯器则多称为玻璃胎珐瑯彩，画珐瑯名称仅用于金属胎者，15但是北京故宫的研究人员也将玻璃胎珐瑯彩器称为玻璃胎画珐瑯。1关于瓷胎珐瑯彩、洋彩与粉彩之名称，上个世纪一直困扰着学界，多数人依循前人的定义，将在瓷胎上以珐瑯釉料彩绘纹饰者称珐瑯彩，初期使用进口釉料，至迟雍正六年（1728）已能自炼珐瑯釉

\vfill
\begin{BookNotes}
\BookNote{13.}{施静菲，〈十八世纪东西交流的见证一清宫画珐瑯工艺在康熙朝的建立〉，《故宫学术季刊》，24卷3期（2007秋季），页59。}
\BookNote{14.}{廖宝秀，〈茶器瑰宝一康熙宜兴胎画珐瑯折枝花卉〉，《故宫文物月刊》，302期，2008年5月，页5，图二，《茶韵茗事一故宫茶话》，页102，图二。廖文中所称《陈设档》乃指称道光十五年七月一日建立的《珐瑯玻璃宜兴磁胎陈设档案》。}
\BookNote{15.}{例如：张临生主编，《故宫鼻烟壶》（台北：故宫，1990)，即如是称。}
\BookNote{16.}{例如夏更起、张荣主编，《故宫鼻烟壶选粹》（北京：故宫，1995)，页109-128。}
\end{BookNotes}

% 原图说　图一　a 清康熙玻璃胎画珐瑯牡丹蓝地胆瓶
% 原图说　b 清康熙玻璃胎画珐瑯牡丹蓝地胆瓶所附木匣　国立故宫博物院藏

% 原书第 428 页（PDF 第 441 页）
\BookPage{428}{441}
\noindent 料。17至于「粉彩瓷器是在康熙五彩基础上，受了珐瑯彩制作工艺的影响而创制的一种釉上彩新品种。」18但无洋彩之名。

近年廖宝秀已基本厘清这个困惑，她比对台北故宫藏品的原名与实物后认为珐瑯彩、粉彩与洋彩都是使用珐瑯釉料彩绘于器胎上者，清宫中没有「粉彩」之名，最早见于清末民初陈浏《匋雅》与许之衡《饮流斋说瓷》等陶瓷专书中。珐瑯彩当然是以珐瑯釉料彩绘，「洋彩」也是珐瑯彩的一种，但是仿用西洋珐瑯画法，故特以「洋」称之，其彩绘地点或完工地点为景德镇御窑场。19清朝内廷无「珐瑯彩」之名，而称「磁（瓷）胎画珐瑯」，其彩绘技法多探用传统中国绘画技法，少数人物探用西洋光影明暗技法，其彩绘地点或完工地点为内廷造办处作坊。此外，两者在题诗、款识与成对表现方式也不同。20

虽然上个世纪所称画珐瑯器大多指称金属胎者，但是朱家滘在其〈清代画珐瑯制器考一《工艺美术史料汇编》之一〉文中所指称的画珐瑯制器，已涵括瓷胎、玻璃胎、金属胎者。

四本乾清宫陈设档册所载录的玻璃胎珐瑯器不多，除了康熙朝的一件「玻璃胎珐瑯牡丹蓝地胆瓶」，就是乾隆朝的三件，即「玻璃胎画珐瑯花卉三喜梅瓶」、「玻璃胎画珐瑯葫芦花插」与「玻璃胎珐瑯八楞瓶」各一件，没有雍正朝玻璃胎珐瑯器。21经检索文物帐，这三件乾隆年间制作而原贮乾清宫库房一端凝殿的玻璃胎珐瑯器，有两件为台北故宫藏品，即「玻璃胎画珐瑯花卉三喜梅瓶」与「玻璃胎珐瑯八楞瓶」。22（图二、图三）

\vfill
\begin{BookNotes}
\BookNote{17.}{冯先铭，《中国陶瓷史（下册）》（台北：胡氏），页387。}
\BookNote{18.}{同前书，页389。}
\BookNote{19.}{廖宝秀，〈乾隆磁胎洋彩综述〉，「二、洋彩、粉彩、珐瑯彩的正名」，页10-15。}
\BookNote{20.}{同前文，「附表：乾隆磁胎洋彩与磁胎画珐瑯区分表」，页31。}
\BookNote{21.}{《珐瑯玻璃宜兴磁胎陈设档案》（道光十五年七月初一日），页1；《珐瑯玻璃宜兴瓷胎陈设档案》(光绪元年十一月初七日)，页1；《金银铜宜兴瓷胎玻璃珐瑯档案》（光绪二年三月初四日)，页46；《明殿现设金银铜宜兴瓷玻璃珐瑯档》（光绪二十八年二月初二日），页86。}
\BookNote{22.}{「玻璃胎画珐瑯花卉三喜梅瓶」的点查号：列六六二，品号：故瓷8528；「玻璃胎珐瑯八楞瓶的点查号：列五六五，品号：故瓷8494；两件的箱号皆是院149。}
\end{BookNotes}

% 原书第 429 页（PDF 第 442 页）
\BookPage{429}{442}
\BookArt{p442-01.jpg}{0.927}{0.314}

检视前述有关玻璃胎、宜兴胎、瓷胎珐瑯器档案，以珐瑯釉料彩绘者大多以「画珐瑯」或「洋彩」称之，仅少数几件只以「珐瑯」称之。依据廖宝秀的研究结论，「洋彩」一词不在本文讨论范围。由于这些陈设档案文物内容多延续抄录，此处仅以光绪二年三月初四日建立的《金银铜宜兴瓷胎玻璃珐瑯档案》为例，仅以「珐瑯」称之的康熙朝瓷胎类有「瓷胎珐瑯菊花白地小瓶」、「瓷胎珐瑯菱花盘」各一件，23雍正朝无此例，在档案倒数第二与第三页的「乾隆款瓷胎」类目中，除了载录了一件「瓷胎珐瑯带耳瓶」24与五份「各样珐瑯炉瓶盒」25外，还载录了一些非珐瑯釉料的瓷器。倒数第三页与「瓷胎珐瑯带耳瓶」同页中，尚有一对「霁红僧帽壶、内壹件火煅」、六份「白瓷红龙杯盘」（分装三匣）、十件「青花白地诗意茶盅」及十件「红花白地诗意茶盅」；到数第二页与「各样珐瑯炉瓶盒」同页者，尚有十对「塡白瓷暗花汤盌」、一对「洋瓷五寸碟」、一对「洋瓷单耳杯」、四十八件「影青釉靶盌」。台北故宫仅藏有一件原贮端凝殿的「霁红僧帽壶」（列四七/故瓷12834/院1736），正合于陈设档中所述有一件「火煅」，可能因此当年南迁时未曾离开紫禁城。此外，「白瓷红龙杯盘」至少有四份（列三九九10/故瓷8520-8523/院109与列三九九12/故瓷8524-8527/院109)，「填白瓷暗

\vfill
\begin{BookNotes}
\BookNote{23.}{《金银铜宜兴瓷胎玻璃珐瑯档案》（光绪二年三月初四日），页7。本文所讨论之康熙朝玻璃胎画珐瑯牡丹蓝地胆瓶亦载录于此页。}
\BookNote{24.}{同前书，页45。}
\BookNote{25.}{同前书，页46。}
\end{BookNotes}

% 原图说　图二　清乾隆玻璃胎画珐瑯花卉三喜梅瓶　国立故宫博物院藏
% 原图说　图三　清乾隆玻璃胎画珐瑯八楞　国立故宫博物院藏

% 原书第 430 页（PDF 第 443 页）
\BookPage{430}{443}
\noindent 花汤盌」至少两对（列三九九13/故瓷17555-17556/院180与列三九九14/故瓷17571-17572/院180)，各一对「洋瓷五寸碟」与「洋瓷单耳杯」也是台北故宫的藏品（列三九九6/故瓷8495-8498/院149)。从这几件实例，可确知它们或是单色釉瓷，或是外国瓷器（洋瓷），或者是白地红彩瓷器。都不是珐瑯釉瓷器，何以列入珐瑯器陈设档案中？至于还有十件「红花白地诗意茶盅」26与六件「瓷胎鸡缸杯」（装一匣）27不排除是珐瑯釉料彩绘，尤其后者，台北故宫也藏有六件原贮端凝殿的乾隆年制瓷胎鸡缸杯（列五五二/故瓷17440-17441/院193、列五五二/故瓷17493-17494/院176、列五五二/故瓷17496/院176、列五五二/故瓷17916/院27)，不排除是在景德镇当地施绘的珐瑯釉彩。

依据乾隆朝《活计档》，28至少在乾隆三年（1738）清高宗已开始为以珐瑯釉料彩绘的各种胎质的器皿配匣装盛，这一年的九月二十九日，他下旨为十五件康熙朝烧制的宜兴胎「烧法瑯」茶具「配匣盛装」，29再入贮乾清宫的库房一端凝殿屋宇中。后来越聚越多，在《活计档》中也一再出现清高宗下令将洋彩或瓷胎珐瑯器配匣后入贮乾清宫「磁器」内，30或「法（珐）瑯（磁）器（皿）」内、31或「磁胎法（珐）瑯器皿」内、32或径称「乾隆款磁珐瑯器皿」内33.。比对目前藏于台北故宫而原贮存于端凝殿的三件玻璃胎珐瑯器，只有一件在乾清宫珐瑯器陈设档中称为「玻璃胎画珐瑯」，其余两件都仅称「玻璃胎珐瑯」。从实物看，它们最大的分别在于乾清宫陈设档中称为「玻璃胎画珐瑯」者乃在不透明白色玻璃（清代内廷名

\vfill
\begin{BookNotes}
\BookNote{26.}{同前书，页45。}
\BookNote{27.}{同前书，页46。}
\BookNote{28.}{中国第一历史档案馆收藏清雍正元年（1723）至宣统三年（1910）的内务府造办处有关活计成做的档案，虽然其中嘉庆前期（四年至十年）似阙遗，然雍正、乾隆两朝档案甚伙，今该馆将其制成微卷，本文所称《活计档》即此。}
\BookNote{29.}{《活计档》，乾隆三年九月二十九日，〈乾清宫〉，即「宜兴烧法瑯盖碗十件、宜兴烧法瑯茶元四件、宜兴烧法瑯茶吊一件」。}
\BookNote{30.}{同前书，乾隆七年四月一日，〈乾清宫〉。}
\BookNote{31.}{例如：同前书，乾隆七年四月六日，〈乾清宫〉；同前书，乾隆九年五月三日，〈如意馆〉；同前书，乾隆九年十二月二十七日，〈灯作〉；同前书，乾隆十年二月十二日，〈乾清宫〉；同前书，乾隆十年二月二十五日，〈乾清宫〉。}
\BookNote{32.}{例如：同前书，乾隆七年四月二十九日，〈乾清宫〉；同前书，乾隆八年十月二十七日，《乾清宫〉；同前书，乾隆九年十二月十八日，〈灯作〉。}
\BookNote{33.}{例如：同前书，乾隆九年六月十八日，〈如意馆〉。}
\end{BookNotes}

% 原书第 431 页（PDF 第 444 页）
\BookPage{431}{444}
\noindent 「涅白」，《活计档》或作「呆白」）胎表以珐瑯釉料彩绘菊花等花卉与彩蝶，器底圈足内有蓝料宋体「乾隆年制」四字款。仅称「玻璃胎珐瑯」者有不同色地，康熙朝的「玻璃胎珐瑯牡丹蓝地胆瓶」，在变形牡丹花与枝叶外，涂满蓝色釉为地；乾隆朝的「玻璃胎珐瑯八楞瓶」因形制较复杂，故彩绘方式也较繁复。八楞瓶的器形又称为「八不正」，器腹好似在四方体上节去八个尖角，腹部有四个正面菱形，其每一边各接一个三角形，虽然每一个正面菱形内皆仅在不透明白色玻璃胎上彩绘五福（蝠）捧团寿字纹饰，但是在每一个团寿字内、每一个三角形内与撇口长颈纹饰之外皆涂满明黄色珐瑯釉色为地。也就是说，这两件仅载录作「玻璃胎珐瑯」器与这件载录作「玻璃胎画珐瑯」器之分别，似乎是在有无色地，而且皆无款。是否这就是道光朝以后内廷当时称为「玻璃胎珐瑯」与称作「玻璃胎画珐瑯」之分别，因仅此三件，尚不能遽下定论。不过，乾隆时期对于这件康熙朝玻璃胎珐瑯胆瓶则名为「玻璃胎画珐瑯牡丹蓝地胆瓶」。（见下文）

这件康熙朝玻璃胎画珐瑯牡丹蓝地胆瓶高12.7公分、口径2.99公分、底径3.95公分、深约11.8公分，乾隆朝的玻璃胎画珐瑯花卉三喜梅瓶高10公分、口径2公分、底径2.2公分，玻璃胎画珐瑯八楞瓶高13.1公分、口径3.5公分、底径3.4公分。器形上八楞瓶最复杂，梅瓶肩部尚饰三个涂金半环系，所以胆瓶形制较简洁。乾隆朝的玻璃胎画珐瑯花卉三喜梅瓶器底圈足内的蓝料宋体「乾隆年制」四字款，当年清室善后委员会点查时已加载录，34另两件无款的玻璃胎画珐瑯器，当年点查时将康熙朝的这件载录为「带盖无款」，35将乾隆朝者载录为「带匣，乾隆年制」。36然而胆瓶不应有盖，当时清室善后委员会出组人员点查时或记载有误。(见下文)

康熙朝玻璃胎画珐瑯牡丹蓝地胆瓶所用的珐瑯釉料有褐紫色（猪肝色）、粉绿色、黄色与蓝色，器腹与颈部皆各彩绘三朵变形褐紫色牡丹花，腹部者大，颈部者小，近足缘则饰三朵紫色团花，花朵内皆加绘粉绿色花瓣；黄色枝叶缠绕于花朵

\vfill
\begin{BookNotes}
\BookNote{34.}{清室善后委员会刊行，《清宫物品点查报告·第二辑》，页211，第一行有「（带破匣，东一号，乾隆年制无圈四字款）」的记载，乃页210最后一行点查号：列六六二条之文物附注。}
\BookNote{35.}{同前书，页186，点查号：列三六73条旁附注:「（带盖无款）」。}
\BookNote{36.}{同前书，页205，点查号:列五六五条旁附注：「（带匣，乾隆年制）」。}
\end{BookNotes}

% 原书第 432 页（PDF 第 445 页）
\BookPage{432}{445}
\noindent 间，叶片上以褐紫色描绘叶脉与阴影；花叶空间塡满蓝色珐瑯釉。前述珐瑯釉料中，粉绿色部位出现不少的鬃眼；色地的蓝色由上而下渐深，显示焙烧时蓝釉向下流聚的现象。

这件玻璃胎画珐瑯牡丹胆瓶附带的木匣，呈上下插盖式，高16.5公分、横长9.8公分、纵长9.64公分，盖面阴刻塡蓝楷书「康熙年制玻璃胎画珐瑯牡丹蓝地胆瓶一件」数字。木匣内有一夹板可卡住瓶身以免玻璃瓶在辖内晃动，但是因此器足将无法碰触匣内底部，或许原应有一木座。若然，则清室善后委员会出组人员点查时的记载或应修订为「带座无款」？（图一b）此木匣应为乾隆时期清高宗下旨配制盛装的，匣盖上的名称为当时的订名。也就是说，虽然这是一件无款玻璃器，但是乾隆时期已明确指出它制作于康熙年间，乾清宫珐瑯器陈设档将它列于「康熙款玻璃胎」类目中，易让人误以为是有款器。两岸故宫所藏清宫文物中经常可见附带木匣、木座的情形，从这件玻璃胎画珐瑯牡丹胆瓶可知，匣盖或木座上所谓「某某年制」或「某某御制」，并不一定表示器身有同字样的款识，除了这件实例外，另一件当时与其收贮于乾清宫的同一大木箱中的一件康熙时期制作的宜兴胎画珐瑯小提梁壶，其匣盖品名阴刻「康熙年制」四字，器腹底部实为「康熙御制」款。37

除此之外，这件康熙朝玻璃胎画珐瑯牡丹蓝地胆瓶的木匣附贴一张破损黄签，其上墨书:「康熙玻璃胎一件」，另附清室善后委员会点查时的签条，左下方墨书:「壹件」。木匣插盖上还朱书「北字一号、一笔、完」数字，匣身上方侧面先朱书再墨写盖之，有楷体「辰一」二字。后者或为此件玻璃胎画珐瑯牡丹蓝地胆瓶收贮在清朝内廷时曾被赋予的编号；至于前者，或即因其收贮在端凝殿北侧小屋，内廷执事人员所赋予的编号与记录。38

康熙朝玻璃胎画珐瑯工艺的最高成就

若以珐瑯釉料施绘于其他胎质表面成器者为画瑯器，则可分为金属胎、陶

\vfill
\begin{BookNotes}
\BookNote{37.}{廖宝秀，茶器瑰宝—康熙宜兴胎画珐瑯折枝花卉〉，《故宫文物月刊》，302期，页5；《茶韵茗事一故宫茶话》，页105。}
\BookNote{38.}{同前注，图九所公布康熙御制宜兴胎画法瑯提梁壶外贮木匣，其插盖面亦有朱书楷体「宜北九号数字。}
\end{BookNotes}

% 原书第 433 页（PDF 第 446 页）
\BookPage{433}{446}
\noindent 瓷胎与玻璃胎三大类。清代康、雍、干三朝画珐瑯器之金属胎有金胎、银胎与铜胎者，以铜胎最常见；陶瓷胎则有宜兴胎与景德镇烧制的瓷胎两类；玻璃胎则仅有不同色泽的玻璃胎之别。以烧制难易言，施绘珐瑯釉料于玻璃胎最为困难，其原因与各类胎质与珐瑯釉料的熔点密切相关。39

黄金的熔点是1063℃，白银的熔点是960℃，纯铜的熔点是1065℃-1083℃，景德镇瓷胎的烧成温度达1200℃，宜兴紫砂胎的烧成温度在1100℃左右。玻璃是非晶体，因成分与添加物的不同，有不同的熔点，1200℃时开始逐渐成为液状，1450℃是其熔融点，但是在温度升高中就先慢慢软化，一般在600℃就有下坠现象，800℃就软化得像麦芽糖般开始流动。珐瑯釉是玻化的物质，在750℃左右就会熔化。相较于其他金属胎与陶瓷胎而言，因为玻璃与珐瑯釉料融程相近，在玻璃胎上施绘珐瑯釉彩相对困难多了。焙烧时温度不足，珐瑯釉料未充分熔解，则呈色不纯正；如果温度过高，则器胎将变形。此外，不同颜色的珐瑯釉料，其融化的温度也不完全相同，在玻璃胎上以珐瑯釉料彩绘纹饰时，需依其熔点高低分别上色，熔点高者先画先烧，反之则后画后烧；一件作品往往需经过三或四次的焙烧方始完成。「当时焙烧的燃料是木炭，又缺乏科学仪器测量温度，只能用眼力和经验判断。所以，为了准确地控制温度以及出炉便捷，只能烧造小件器物。即便如此，失败的比例也很高。」40

经过中外学者多年的研究，41今人已能比较清晰的明了清初康、雍、干三朝内

\vfill
\begin{BookNotes}
\BookNote{39.}{夏更起，〈玻璃胎画珐瑯考析〉，《故宫博物院院刊》，2003年3期，页18，「2.制造之艰难」一节作者首从熔点论证，认为铜胎、瓷胎画珐瑯器的「胎体熔点均比珐瑯釉料高出数百度，彩绘后的焙烧技术较易掌握。」}
\BookNote{40.}{同前注。}
\BookNote{41.}{详细相关资料可先参阅林姝，〈雍正时期玻璃制品与朝政的关系〉，《故宫博物院院刊》，2008年5期，页45，注释1。此外尚有：Peter Hardie, "Some Decorative Techniques Found in Later Chinese Glass," Ransactions of the Oriental Ceramic Society, 1979-1980, pp. 47-66. Peter Hardie , "Glass in China: LateMing to Early Qing,"Ransactions of the OrientalCeramic Society, 1990-1991, pp.9-28. Emily Byrne Curtis (edited), Pure brightness shines everywhere : the glass of China, England: Ashgate, 2004. Emily Byrne Curtis, Glass Exchange between Europe and China, 1550-1800, England: Ashgate, 2008.施静菲，〈十八世纪东西交流的见证一清宫画珐瑯工艺在康熙朝的建立〉。许晓东，〈康熙、雍正时期宫廷与地方画珐瑯技术的互动〉，收于故宫博物院、柏林马普学会科学}
\end{BookNotes}

% 原书第 434 页（PDF 第 447 页）
\BookPage{434}{447}
\noindent 廷烧制玻璃器的发展梗概，作为清代玻璃制作工艺的开端与奠定基础时期的康熙朝，「已有单色玻璃、套玻璃、珐瑯彩玻璃（即玻璃胎画珐瑯）、刻花玻璃、洒金玻璃等几个品种。」42近人也更针对康熙朝与雍正朝的金属胎、玻璃胎与瓷胎画珐瑯器，论证康熙时期画瑯工艺所显示的中西交流，43以及康、雍二朝宫廷与地方画珐瑯技术的互动。44

如今我们可知，清圣祖甚重视玻璃器与珐瑯器，内廷造办处玻璃厂与珐瑯作（处）内除了有汉满匠役外，45外国传教士也参与研制。自康熙三十五年（1696）玻璃厂设置于当时北京皇城西安门内北堂的蚕池口附近，46至康熙五十七年（1718）武英殿珐瑯作改归养心殿，称珐瑯处，47规模与编制扩大，匠役与传教士已能充分掌握玻璃胎画珐瑯工艺技术，当时烧制的成品不但供皇帝自己玩赏，也作为赏赐大臣的礼物，以示荣宠。48此外，皇帝也曾将玻璃胎画珐瑯器皿赏赐西洋人

\vfill
\begin{BookNotes}
\BookNote{}{史所编，《宫廷与地方:十七至十八世纪的技术交流》（北京:紫禁城，2010)，页277-335。等等。}
\BookNote{42.}{张荣，〈清康熙朝御制玻璃〉，收于《明清论丛·第二辑》（北京:紫禁城，2001），页539-543。}
\BookNote{43.}{施静菲，〈十八世纪东西交流的见证一清宫画珐瑯工艺在康熙朝的建立〉，页45-94。}
\BookNote{44.}{许晓东，<康熙、雍正时期宫廷与地方画法瑯技术的互动〉。}
\BookNote{45.}{《活计档》，雍正三年十一月，〈玉作〉，载：广东匠人程向贵制作的一套十二件「雨过天晴刻花套杯」乃于康熙四十八年正月出八日呈进的；另一套十二件「雨过天晴素套杯」是「家里学手匠人周俊」做的，于雍正二年十月二十八日呈进。经查：「广东匠人程向贵、周俊二人于康熙五十四年四月交巡抚送去讫。}
\BookNote{46.}{林业强，〈清宫玻璃厂三题〉，《光凝秋水—清宫造办处玻璃器》（北京：紫禁城，2005），页28}
\BookNote{43。}{EmilyByrneCurtis，米辰峰译，〈清朝的玻璃制造与耶稣会士在蚕池口的作坊〉，《故宫博物院院刊》，2003年1期(总105)，页62-71。原文"QingGlassmaking: the Jesuit Workshop onCanchikou"刊于《光凝秋水—清宫造办处玻璃器》，页94-102，书中亦将该文中译后登出，但篇名为〈耶稣会士与清朝的玻璃制造〉，页84-93。民国初期，蚕池口就只改为国务院、总统府等，三年代为河北省政府所在地。见林业强，清宫玻璃厂三题〉，页31，最后一行。}
\BookNote{47.}{康熙十九年（1680）武英殿设造办处，但是养心殿仍有皇帝最重视的匠役在其中工作，康熙四十七年（1708）清圣祖下旨将养心殿的匠役都移至造办处，可是康熙五十七年又将武英殿珐瑯作改归养心殿扩大编制成珐瑯处，需至雍正朝方再移出至武英殿。}
\BookNote{48.}{国立故宫博物院编，《宫中档康熙朝奏折·第六辑》（台北：故宫，1976)，页603，时任广西巡抚兵部左侍郎兼都察院右副都御史陈元龙于康熙五十五年九月十一日奏谢清圣祖恩赐「御制法琅五彩红玻璃鼻烟壶一、八角盒砚一、水丞一、圆香盒一。…恩赐法琅宝器四种，…」。其中仅「法琅五彩红玻璃鼻烟壶为康熙御制玻璃胎画珐瑯器皿。EmilyByrne Curtis 文中(Qing Glassmaking: theJesuit Workshop on Canchikou)引用《宫中档康熙朝奏}
\end{BookNotes}

% 原书第 435 页（PDF 第 448 页）
\BookPage{435}{448}
\noindent 士，以展现内廷研制这种工艺的成就。49

虽然如此，以目前所能获得的资料以及传世实物数量看来，当时内廷匠役或传教士较能掌控金属胎与陶瓷胎画珐瑯器皿的烧制，烧制的成品较多，尺寸也较大。在乾清宫珐瑯器陈设档案中也显示出这样的现象，以道光十五年建立的七月初一日建立的《珐瑯玻璃宜兴磁胎陈设档案》为例，当时贮存于端凝殿小屋中金属胎之外的珐瑯器，康熙朝宜兴胎画珐瑯有二十件，瓷胎者有十六件，玻璃胎仅一件。50由同时建立的《珐瑯金银铜胎陈设档案》可知，当时贮存于端凝殿小屋的金属胎珐瑯器中，康熙朝金胎珐瑯器有四件，径称「画珐瑯」者仅一件；51康熙朝铜胎珐瑯器除了有画瑯器，尚有广珐瑯器、西洋珐瑯器，以及不少径称「珐瑯」者，称为「铜胎画瑯」者也有二十一件。52

康熙朝御制的玻璃胎画珐瑯器当然不仅只有这件贮存于乾清宫的玻璃胎画珐瑯牡丹蓝地胆瓶，例如赏赐给大臣陈元龙的就有一件画珐瑯五彩红玻璃鼻烟壶；53今人香港李景勋也藏有一件高仅3.6公分、口径5.1公分的画珐瑯花卉小杯，其器外壁四开光内彩绘四季花卉，分别是牡丹、荷花、菊花与月季，开光外上下缘绘饰半团形图案式花卉纹，中间塡以红彩，器底内凹，正中以蓝料描双方框，内书宋体「康熙御制」款；54惜于上世纪在展览中受损，幸仍可黏合成器。（图四）此外，欧

\vfill
\begin{BookNotes}
\BookNote{}{折》（台北：故宫，1976)，第六辑，页603广西巡抚兵部左侍郎兼都察院右副都御史陈元龙的谢恩折，亦仅认为鼻烟壶是玻璃胎珐瑯器皿，但是译者却混为「四件画珐瑯玻璃器皿」。又，EmilyByrneCurtis文中引用《宫中档康熙朝奏折》第七辑，页357-361，广西提督总兵官左都督左世永的康熙五十七年六月十三日谢恩折，认为清圣祖颇以内务府造办处珐瑯匠役技艺的精进为荣，因为他曾将内廷制作的珐瑯器皿赏赐大臣；译者则误译作清圣祖将「画珐瑯玻璃器皿」赏赐重臣。而档案原文为御制法瑯水盛壹个、嵌花石紫石盒砚一方、双联荷包壹个，又镶荷包内法瑯鼻烟壶壹个，又黄荷包内火镰包壹个」，所以此次的赏赐有两件珐瑯器皿，但不确定是否是玻璃胎画珐瑯器，亦不排除乃铜胎掐丝珐瑯器。}
\BookNote{49.}{Emily Byrne Curtis, "Qing Glassmaking: the Jesuit Workshop on Canchikou," p. 95, "……Kangxi sent to de Tournon, an enameled glass snuff bottle of His Majesty's own use."}
\BookNote{50.}{《珐瑯玻璃宜兴磁胎陈设档案》（道光十五年七月初一日)，页1-3。}
\BookNote{51.}{《珐瑯金银铜胎陈设档案》（道光十五年七月初一日），页1。}
\BookNote{52.}{《珐瑯金银铜胎陈设档案》(道光十五年七月初一日)，页3-7。}
\BookNote{53.}{国立故宫博物院编，《宫中档康熙朝奏折·第六辑》，页603。}
\BookNote{54.}{许建勋、林业强合著，黎淑仪编，《虹影瑶辉—李景勋藏清代玻璃》（香港：中大，2000)，页364-365，No.144°}
\end{BookNotes}

% 原书第 436 页（PDF 第 449 页）
\BookPage{436}{449}
\BookArt{p449-01.jpg}{0.990}{0.418}

\noindent 洲荷兰的海牙博物馆(HaagsGemeentemuseum)则藏有一件康熙款玻璃胎花卉黄地盒，圈足内正中亦以蓝料描双方框，内书宋体「康熙御制」款识，其纹饰为西番花，55在台北故宫珍藏的康熙朝瓷胎画珐瑯碗中可觅得类似黄地花卉纹饰与蓝料款识；56至于器形，也有可资比对的相类康熙朝瓷胎画珐瑯盒。57从已知的这三件康熙朝玻璃胎画珐瑯器，以台北故宫所藏这件胆瓶的尺寸最大，当然烧制工艺更高；从现存六本乾清宫库收珐瑯器陈设档案可知，它是清朝唯一一件乾清宫的库房一端凝殿收贮的康熙朝玻璃胎珐瑯器，或即意味着清帝认为它是康熙朝玻璃胎画珐瑯制作工艺的最高成就！

\vfill
\begin{BookNotes}
\BookNote{55.}{Peter Hardie, "Glass inChina: Questions andSome Answers," p.10, fig. 7 (黑白图). Emily Byrne Curtis, Glass Exchange between Europe and China, 1550-1800, p. 55, fig. 6.9\&fig. 6.10 (器为彩色图，款识为黑白图）}
\BookNote{56.}{蔡和璧编，《清宫中珐瑯彩瓷特展》（台北:故宫，1992)，页39-40，图3与图4。}
\BookNote{57.}{同前书，页43，图7，蔡合璧认为这件紫红地番莲小盒就是道光十五年七月一日建立的《珐瑯玻璃宜兴磁胎陈设档案》，所列的「磁胎画珐瑯西番花红地盒」，高1.9公分、通盖高3.1公分、口径5.3公分、足径3.5公分。经查原陈设档案，「康熙款磁胎」类目中共有十种二十二件康熙磁胎画珐瑯器，最后一件为唯一的瓷盒，乃「磁胎画珐瑯五彩西番花红地盒」。}
\end{BookNotes}

% 原图说　图四　清康熙玻璃胎画珐瑯花卉小杯　香港李景勋藏

% 图表 1 的文字辨识草稿，行列按原图查看。
% 御康
% 制配
% 图四清康熙玻璃胎画珐瑯花卉小杯香港李景勋藏

% 原书第 437 页（PDF 第 450 页）
\BookPage{437}{450}
\BookArt{p450-01.jpg}{0.990}{0.830}


% 原图说　附表　光绪二年三月初四日《金银铜宜兴瓷胎玻璃珐瑯档案》与光绪二十八年二月初二日《明殿现设金银铜宜兴瓷玻璃珐瑯档》对照表

% 图表 1 的文字辨识草稿，行列按原图查看。
% 附表光绪二年三月初四日《金银铜宜兴瓷胎玻璃珐瑯档案》与光绪二十八年二月初二日《明殿现设金银
% 铜宜兴瓷玻璃珐瑯档》对照表
% 页行
% 《金银铜宜兴瓷胎玻璃珐瑯档案》
% 《明殿现设金银铜宜兴瓷玻璃珐瑯档》
% 数数
% 光绪二年三月初四日奴才范常禄等陆
% 光绪二十八年二月初二日总管李长喜等陆
% 1-2
% 42 1-2
% 续查得乾清宫库收金银铜胎宜兴瓷胎
% 续查得乾清宫库收三朝珐瑯
% 玻璃珐瑯
% 铜胎（旁边贴黄签:「西洋」）珐瑯
% 45
% 铜胎西洋珐瑯杯盘壹分
% 杯盘壹分
% 铜胎（「广」用黄签黏遮）珐瑯小盒
% 45
% 铜胎珐瑯小盒壹件、缺瑯
% 壹件、缺瑯
% 铜胎（「广」用黄签黏遮）珐瑯小盒
% 45
% 铜胎珐瑯小盒壹件
% 壹件
% 铜胎珐瑯筒子炉贰件、（下方黏贴黄
% 签）光绪二十一年又五月二十二日内
% 49
% 铜胎珐瑯筒子炉壹件
% 殿总管谢双喜传上要去铜胎珐瑯筒
% 子炉一件赏达赖喇嘛用
% 铜胎珐瑯甘露壶壹件、（下方黏贴黄
% 无此条，将下一页第一行的内容顺移至此
% 签）光绪二十一年又五月二十二日内
% 50
% 10
% 殿总管谢双喜传上要去铜胎珐瑯甘
% 页第十行（以下至52页皆同）
% 露壶一件赏达赖喇嘛用
% 瓷胎画珐瑯久安长治白地大盌壹对
% 11
% （「久」字乃用黄色签另书写后黏贴
% 51
% 瓷胎画珐瑯久安长治白地大盌壹对
% 上，下方原字无法辨识)
% 瓷胎画珐瑯水墨牡丹白地宫盌壹对、
% 12
% 52
% 瓷胎画珐瑯水墨牡丹白地宫盌壹对、有缺
% 有缺（「牡丹」字乃用黄色签另书写
% 后黏贴上，下方原字无法辨识）
% 瓷胎画珐瑯（旁边贴黄签:「五
% 53
% 13
% 瓷胎画珐瑯五色西番花红地茶盌壹对
% 色」）西番花红地茶盌壹对
% 瓷胎画珐瑯仙芝祝寿黄地茶盅壹件
% 瓷胎画珐瑯仙芝祝寿黄地茶盅壹件（此页
% (「黄」字乃用白色签另书写后
% 13
% 53
% 10
% 共10行，故末行文字相同）
% 黏贴上，下方原「茶」字。此页
% 共11行）
% 瓷胎画砝瑯（「青」用黄签黏遮）山
% 57
% 17
% 瓷胎画珐瑯山水六寸盘壹对
% 11
% 水六寸盘壹对
% 瓷胎画珐瑯红梅花黄地六寸碟壹对
% 12
% 58
% 瓷胎画珐瑯红梅花黄地六寸碟壹对
% 18
% (「红」字乃用白色签另书写后黏贴
% 上，下方原字无法辨识）
% 金胎珐瑯金瓯辑瑞壹对、计一匣、重贰拾
% 金胎珐瑯金瓯辑瑞壹对、计一匣、重
% 两伍钱壹件、重拾玖两伍钱壹件、（下方
% 13
% 20
% 60
% 贰拾两伍钱壹件、重拾玖两伍钱壹件
% 黏贴黄签）光绪三十年二月三十日海总管
% 阮进寿传上要去金胎金瓯辑瑞壹对、完
% 金胎珐瑯爹斗贰对、…(下方黏贴黄签)
% 62
% 22
% 金胎珐瑯爹斗贰对、…
% 光绪三十年二月三十日海总管阮进寿传
% 14
% 上要去金胎珐瑯金爹斗壹对、完

% 原书第 438 页（PDF 第 451 页）
\BookPage{438}{451}
\BookArt{p451-01.jpg}{0.990}{0.830}

\vfill
\begin{BookNotes}
\BookNote{58.}{此黄签文字疑点有二，1、农历二月无三十日，2、随「坐」应是随「座」。}
\end{BookNotes}

% 图表 1 的文字辨识草稿，行列按原图查看。
% 15
% 22
% 10
% 62
% 光绪三十年二月三十日海总管阮进寿传
% 金胎珐瑯盖盌壹对、
% 上要去金胎珐瑯盖盘壹对、完
% 金胎珐瑯盖罐壹对、…(下方黏贴黄签)
% 63
% 16
% 23
% 光绪三十年二月三十日海总管阮进寿传
% 金胎珐瑯盖罐壹对、…
% 上要去金胎珐瑯盖罐壹对、完
% 铜胎广珐瑯西番花黄地炉瓶香盒壹
% 铜胎广珐瑯西番花黄地炉瓶香盒壹分、年
% 65
% 17
% 25
% 分、年节用（「广」字乃用黄签另书
% 节用
% 写后黏贴上，下方原字无法辨识）
% 铜胎掐丝珐瑯异兽双环扁瓶壹件
% 26
% 66
% 铜胎掐丝珐瑯异兽双环扁瓶壹件
% 18
% 10
% (「双环」字乃用白色签另书写后黏
% 10
% 贴上，下方原字无法辨识）
% 铜胎掐丝瑯盘大小拾壹件、（下方
% 黏贴黄签）光绪五年闰二月二十二日
% 72
% 铜胎掐丝珐瑯盘大小玖件
% 19 32
% 内殿总管李双喜要去铜珐瑯盘壹对赏
% 达赖喇嘛呼毕勒罕用
% 瓷胎画珐瑯蝴蝶红地四寸碟壹对
% 瓷胎画珐瑯蝴蝶红地四寸碟壹对
% (「画珐」与「对」字乃用白色签另
% 77 10
% 20
% 37
% 10
% 书写后黏贴上，下方原字无法辨识）
% 「对」字乃用白色签另书写后黏贴
% 上，下方原字无法辨识）（「对」字
% 「对」字乃用白色签另书写后黏贴上，下
% 80
% 21 40
% 方原字无法辨识）
% 乃用白色签另书写后黏贴上，下方原
% 字无法辨识）
% 霁红僧帽壶壹对、内壹件火煅
% 22 44
% 霁红僧帽壶壹对、内壹件火煅
% (「壶」字乃用白色签另书写后黏贴
% 上，下方原字无法辨识）
% 瓷胎鸡缸杯陆件、计壹匣（「缸」字
% 86
% 瓷胎鸡缸杯陆件、计壹匣
% 23
% 46
% 乃用黄签另书写后黏贴上，下方原字
% 无法辨识）
% 填白瓷暗花汤盌对、计拾匣、内有缺玺
% 填白瓷暗花汤盘拾对、计拾匣、内有
% 86
% （「对」字上一字用白色签黏贴遮蔽，依
% 24 46
% 残余笔画言，似「拾」字）
% 洋瓷百鹿壹对、紫檀木座、（下方黏贴
% 光绪八年十二月十六日
% 黄签）光绪三十年二月三十日海总管阮进
% 87
% 25
% 47 1-3
% 长春宫首领杨得清交洋瓷百鹿壹
% 寿传上要去铜胎洋瓷百鹿尊一对、完、
% 对、紫檀木座
% 随坐58
% 光绪十三年三月十二日奏事处首领郭
% 双喜交
% 蓝洋瓷通子瓶壹对无底、（下方黏贴
% 47 4-5
% 26
% 黄签）光绪二十四年闰三月二十七日
% 御前代班春福传上要去蓝洋瓷筩子
% 瓶一对无底、完
